\section{Alcance}

El presente trabajo está enfocado en el diseño e implementación de un prototipo web educativo orientado a apoyar la comprensión de los conceptos fundamentales de interpretación y compilación de lenguajes de programación en la asignatura FLP del programa de Ingeniería de Sistemas de la Universidad del Valle sede Tuluá.

Desde el punto de vista técnico, el prototipo cubre tres funcionalidades principales: (1) la definición y edición de gramáticas con sintaxis configurable por el usuario; (2) la visualización de árboles de sintaxis abstracta (AST) generados a partir de dichas gramáticas; y (3) la representación de los cambios en el ambiente durante la evaluación de programas escritos bajo la gramática definida. Estas funcionalidades están orientadas específicamente a los conceptos que, conforme al diagnóstico presentado en el Capítulo 3, presentan mayor dificultad de comprensión en la asignatura.

El alcance del prototipo se limita a funcionalidades demostrativas con fines pedagógicos. No se pretende desarrollar un lenguaje de propósito general ni implementar mecanismos avanzados de ejecución tales como depuración paso a paso, compilación a código máquina, optimización de rendimiento o autenticación y persistencia de usuarios. El sistema tampoco contempla despliegue en infraestructura institucional de uso permanente, puesto que su operación se restringe a un entorno de pruebas durante la fase de evaluación.

La evaluación del prototipo se llevará a cabo con un grupo de entre 10 y 20 estudiantes matriculados en la asignatura durante el período de desarrollo del trabajo, junto con el docente responsable. Dada la naturaleza y el tiempo disponible del proyecto, la evaluación se restringe a aspectos de funcionalidad técnica y usabilidad, mediante la aplicación del cuestionario estandarizado System Usability Scale (SUS) y un cuestionario de utilidad pedagógica de elaboración propia, dejando por fuera evaluaciones de impacto en el aprendizaje a largo plazo, pues estos requieren condiciones experimentales que exceden el alcance de un trabajo de grado. 

Finalmente, el prototipo no pretende reemplazar los recursos actuales de la asignatura, en particular lo referente a Racket y EOPL, sino complementarlos, ofreciendo un punto de entrada más accesible y visual a los conceptos que los estudiantes identifican como los de mayor dificultad.